\section{Data layout and environments}

\subsection{Where data lives}

\begin{description}
  \item[out/sweeps/] one directory per sweep session: a
    \cmd{config.json}, one trace archive per completed azimuth, and
    fit results. Resumable and analysable while incomplete.
  \item[out/tesscache/, out/observables/] the tessellation cache and
    observable matrices, regenerated on demand.
  \item[sim/analysis/*/ (stored npz)] each analysis theme stores its
    own results beside its code.
  \item[sim/ref/] the clean reference traces behind the paper's
    figures.
\end{description}

None of \cmd{out/} is version-controlled; it is deposited with the
paper. \cmd{DATA\_MANIFEST.md} records the mapping. The large
finite-element meshes are not versioned; the meshing scripts rebuild
them if rerun.

\subsection{Environments}

Three usage tiers, in increasing weight:

\begin{enumerate}
  \item \textbf{Analysis only} (no GPU): the pinned numpy and scipy
    of \cmd{requirements.txt}. Every stored result reproduces on
    CPU.
  \item \textbf{GUI}: the pinned set plus streamlit and plotly
    (matplotlib is already required for figures). The launcher's
    private runtime installs exactly this.
  \item \textbf{New solver runs}: additionally CUDA and cupy. Only
    needed to regenerate sweeps or reference traces.
\end{enumerate}

\subsection{The manual}

This document builds from \cmd{docs/} with \cmd{build.ps1} (XeLaTeX,
two passes). Each chapter is a numbered folder and each section a
separate file, so updates touch a single small file.
